\section*{Capítulo \ref{ch:b1:stress-strain-rheology} --- Esfuerzo, deformación y reología}

\begin{solution}{exo:b1:stress-strain-rheology:1}
(a) $\sigma_v = 2700\times9.81\times10\,000 = \qty{2.65e8}{Pa} = \qty{264.9}{MPa}$.
(b) $P = 1000\times9.81\times10\,000 = \qty{98.1}{MPa}$. (c) $\sigma_v' =
264.9 - 98.1 = \qty{166.8}{MPa}$, y $\lambda = 98.1/264.9 = 0.370$, que es
$1000/2700$.
\end{solution}

\begin{solution}{exo:b1:stress-strain-rheology:2}
(a) $e = (13-10)/10 = 0.30$ y $S = 1.30$. (b) $e = (7-10)/10 = -0.30$, un
acortamiento de \qty{30}{\percent}, y $S = 0.70$. (c) El ángulo se ha cerrado en
$\psi = \ang{30}$, y $\gamma = \tan\ang{30} = 0.577$.
\end{solution}

\begin{solution}{exo:b1:stress-strain-rheology:3}
(a) Centro \qty{70}{MPa}, radio \qty{30}{MPa}. (b) El esfuerzo de cizalla máximo
es el radio, \qty{30}{MPa}, en los planos a \ang{45} de $\sigma_1$ y
$\sigma_3$, que llevan el \omterm{def:b1:stress-strain-rheology:stress}{esfuerzo normal} del centro, \qty{70}{MPa}. (c)
$2\theta = \ang{60}$: $\sigma_n = 70 + 30\cos\ang{60} = \qty{85}{MPa}$ y
$\tau = 30\sin\ang{60} = \qty{26.0}{MPa}$.
\end{solution}

\begin{solution}{exo:b1:stress-strain-rheology:4}
(a) $\ang{45} - \ang{30}/2 = \ang{30}$ respecto de $\sigma_1$, a ambos lados de
él, así que los dos planos forman \ang{60} entre sí. (b) Con $c' = 20 + r$, la
condición de \cref{thm:b1:stress-strain-rheology:failure-angle} se escribe
$r = C\cos\phi + (20 + r)\sin\phi$, así que
\[
  r = \frac{C\cos\phi + 20\sin\phi}{1 - \sin\phi} = \frac{8.66 + 10}{0.5}
    = \qty{37.3}{MPa},
\]
y $\sigma_1' = 20 + 2r = \qty{94.6}{MPa}$.
\end{solution}

\begin{solution}{exo:b1:stress-strain-rheology:5}
(a) $c = 40$ y $R = \sqrt{20^2 + 15^2} = \qty{25}{MPa}$: $\sigma_1 =
\qty{65}{MPa}$ y $\sigma_3 = \qty{15}{MPa}$; $\tan2\theta_p = 30/40 = 0.75$,
$2\theta_p = \ang{36.87}$, $\theta_p = \ang{18.43}$. (b) Con este ángulo
$\cos^2\theta = 0.9$, $\sin^2\theta = 0.1$ y $\sin\theta\cos\theta = 0.3$:
$\sigma_n = 60\times0.9 + 2\times15\times0.3 + 20\times0.1 = 54 + 9 + 2 =
\qty{65}{MPa}$, y $\tau = (20-60)\times0.3 + 15\times(0.9-0.1) = -12 + 12 =
0$. (c) $R = \qty{25}{MPa}$, en los planos a \ang{45} de las direcciones
principales.
\end{solution}

\begin{solution}{exo:b1:stress-strain-rheology:6}
(a) $S_{1,2} = \sqrt{1 + 1} \pm 1$: $S_1 = 2.414$ y $S_2 = 0.414$;
$\tan2\theta' = 2/2 = 1$, $\theta' = \ang{22.5}$. (b) $S_1S_2 =
(\sqrt2 + 1)(\sqrt2 - 1) = 1$, que es $\det F$: no hay cambio de área. (c) La
línea según $(1, 1)$ pasa a $F(1, 1) = (1 + 2, 1) = (3, 1)$, a
$\arctan(1/3) = \ang{18.4}$ de la dirección de cizalla, con un estiramiento
$\sqrt{10}/\sqrt2 = 2.24$. Ha girado más allá del eje largo de la elipse, a
\ang{22.5}, hacia el plano de cizalla; su estiramiento es algo menor que
$S_1$, como debe ser, ya que ninguna línea se estira más que $S_1$.
\end{solution}

\begin{solution}{exo:b1:stress-strain-rheology:7}
(a) $S_1/S_2 = 4$ y $S_1S_2 = 1$ dan $S_1 = 2$ y $S_2 = 0.5$; según
\cref{prop:b1:stress-strain-rheology:simple-shear}, $\gamma = S_1 - S_2 =
1.5$. (b) $\tan2\theta' = 2/1.5 = 1.333$, $2\theta' = \ang{53.1}$: los ejes
largos deben situarse a \ang{26.6} de las paredes. (c) $1.5\times200 = \qty{300}{m}$.
\end{solution}

\begin{solution}{exo:b1:stress-strain-rheology:8}
$\sigma_v = 2700\times9.81\times3000 = \qty{79.46}{MPa}$. (a) $P = \qty{29.43}{MPa}$,
$\sigma_n' = 70 - 29.43 = \qty{40.57}{MPa}$, y la \omterm{prop:b1:stress-strain-rheology:byerlee}{ley de Byerlee} permite
$0.85\times40.57 = \qty{34.5}{MPa}$, más que los \qty{30}{MPa} aplicados: la
falla aguanta. (b) $P = 0.6\times79.46 = \qty{47.68}{MPa}$, $\sigma_n' =
\qty{22.32}{MPa}$, y la fricción da solo \qty{19.0}{MPa}: la falla
desliza. (c) $30 = 0.85\,(70 - 79.46\,\lambda)$ da $\lambda = (70 -
35.29)/79.46 = 0.437$.
\end{solution}

\begin{solution}{exo:b1:stress-strain-rheology:9}
(a) El esfuerzo varía como $\dot\varepsilon^{1/4}$: $10^{1/4} = 1.78$. (b) El
cociente de las resistencias es
\[
  \exp\Bigl[\frac{E}{nR}\Bigl(\frac{1}{673.15} - \frac{1}{773.15}\Bigr)\Bigr]
  = \exp\bigl(6705\times\num{1.921e-4}\bigr) = \exp(1.288) = 3.63 .
\]
(c) Un cambio de diez veces en la velocidad cambia la resistencia de fluencia
en un factor 1.78, que la temperatura produce en solo $100\times\ln1.78/\ln3.63 \approx
\qty{45}{\degreeCelsius}$; con una geoterma de unos \qty{20}{\degreeCelsius}
por kilómetro, eso desplaza la transición unos \qty{2}{km}. Un cambio de
geoterma desplaza las isotermas, y con ellas la transición, mucho más.
\end{solution}

\begin{solution}{exo:b1:stress-strain-rheology:10}
(a) Con $C = 0$, una falla desliza cuando $r = \mu c'/\sqrt{1+\mu^2}$
(\cref{thm:b1:stress-strain-rheology:failure-angle}). Aquí $c' = S =
\rho g z(1-\lambda)$, así que $\sigma_1 - \sigma_3 = 2r = 2\mu S/\sqrt{1+\mu^2}$.
(b) $S$ crece $2700\times9.81\times1000\times(1 - 1000/2700) =
\qty{16.68}{MPa}$ por kilómetro, y $2\times0.85/1.3124 = 1.295$: la pendiente es
\qty{21.6}{MPa} por kilómetro, entre \num{13.1} (fallas normales) y
\qty{61.3}{MPa} por kilómetro (fallas inversas). (c) $2\mu/(k+\mu) = 0.786$,
$2\mu/k = 1.295$ y $2\mu/(k-\mu) = 3.68$, con $k = 1.3124$: el \omterm{def:b1:stress-strain-rheology:principal}{esfuerzo
diferencial} en el deslizamiento es 0.79, 1.30 y 3.68 veces el esfuerzo
vertical efectivo.
\end{solution}

\begin{solution}{exo:b1:stress-strain-rheology:11}
(a) $t = 60\times\num{e6}\times\num{3.156e7} = \qty{1.89e15}{s}$ y
$2\sqrt{\kappa t} = 2\sqrt{\num{1.89e9}} = \qty{87.0}{km}$. (b)
$\operatorname{erf}x = 600/1350 = 0.444$ para $x = 0.417$: $z = 0.417\times87.0
= \qty{36.3}{km}$; $\operatorname{erf}x = 800/1350 = 0.593$ para $x = 0.586$:
$z = \qty{51.0}{km}$. (c) Los terremotos deberían alcanzar \qtyrange{36}{51}{km};
la envolvente sitúa la base de la capa frágil a \qty{38.2}{km}, dentro de ese
intervalo y cerca de su parte superior, a una temperatura, \qty{641}{\degreeCelsius}, que
está en el intervalo observado.
\end{solution}

\begin{solution}{exo:b1:stress-strain-rheology:12}
(a) $r = 45$ y $c' = \qty{65}{MPa}$; la distancia del centro a la recta
$\tau = 0.85\sigma_n'$ es $0.85\times65/1.3124 = \qty{42.1}{MPa}$, menor que
el radio: el círculo corta la recta, y los planos del arco situado por encima
de ella deslizan. (b) Con $2\theta$ medido desde el extremo de $\sigma_1$, el
deslizamiento exige $45\sin2\theta \ge 0.85\,(65 + 45\cos2\theta)$, es decir
$45\sin2\theta - 38.25\cos2\theta \ge 55.25$, o bien
$59.06\sin(2\theta - \ang{40.36}) \ge 55.25$, donde $\ang{40.36} = \phi$.
Así que $\sin(2\theta - \phi) \ge 0.9355$, $2\theta - \phi$ entre \ang{69.31}
y \ang{110.69}, $2\theta$ entre \ang{109.67} y \ang{151.06}: las normales
están a entre \ang{54.8} y \ang{75.5} de $\sigma_1$, y las fracturas que
deslizan, a entre \ang{14.5} y \ang{35.2} de $\sigma_1$. El centro del intervalo,
\ang{24.8}, es $\ang{45} - \phi/2$, el plano de
\cref{thm:b1:stress-strain-rheology:failure-angle}. (c) Solo ese plano desliza
cuando el círculo toca justo la recta: $r = \mu(20 + r)/k$, $r = 20\mu/(k -
\mu) = 17/0.4624 = \qty{36.8}{MPa}$, y $\sigma_1' = 20 + 2r =
\qty{93.5}{MPa}$.
\end{solution}

\begin{solution}{pb:b1:stress-strain-rheology:1}
\textbf{1.} $\sigma_v = 2700\times9.81\times5000 = \qty{132.44}{MPa}$.

\textbf{2.} $P = 1000\times9.81\times5000 = \qty{49.05}{MPa}$, y
$\lambda = 49.05/132.44 = 0.370$.

\textbf{3.} $\sigma_1 = \sigma_v = \qty{132.44}{MPa}$, vertical, y
$\sigma_3 = \qty{75.00}{MPa}$, horizontal este--oeste: la \omterm{def:b1:layered-earth:layers}{corteza} se estira
de este a oeste bajo su propio peso, el estado que impulsa las fallas
normales, como esta.

\textbf{4.} $\sigma_1 - \sigma_3 = \qty{57.44}{MPa}$; centro
$(132.44 + 75.00)/2 = \qty{103.72}{MPa}$; radio \qty{28.72}{MPa}.

\textbf{5.} $\sigma_1' = 132.44 - 49.05 = \qty{83.39}{MPa}$, $\sigma_3' =
75.00 - 49.05 = \qty{25.95}{MPa}$, centro $103.72 - 49.05 = \qty{54.67}{MPa}$.

\textbf{6.} $\sigma_{xx} = \qty{75.00}{MPa}$, $\sigma_{yy} = \qty{132.44}{MPa}$
y $\sigma_{xy} = 0$: los ejes son las direcciones principales, una de ellas
vertical, así que los planos perpendiculares a ellas no llevan esfuerzo de cizalla.

\textbf{7.} Según su \omterm{def:b1:geological-maps:attitude}{buzamiento}, la falla sigue el vector unitario
$(\cos\ang{60}, -\sin\ang{60})$, hacia el este y hacia abajo; su normal hacia arriba
es $\vec n = (\sin\ang{60}, \cos\ang{60}) = (0.866, 0.500)$, a \ang{30} de
$x$. Su ángulo con $\sigma_1$, según $y$, tiene coseno
0.500: $\theta = \ang{60}$.

\textbf{8.} Con \ang{30} desde $x$: $\sigma_n = 75.00\times0.75 + 132.44\times0.25
= 56.25 + 33.11 = \qty{89.36}{MPa}$, y $\tau = (132.44 - 75.00)\times
\sin\ang{30}\cos\ang{30} = 57.44\times0.4330 = \qty{24.87}{MPa}$.

\textbf{9.} A $2\theta = \ang{120}$ desde el extremo de $\sigma_1$:
$\sigma_n = 103.72 + 28.72\cos\ang{120} = 103.72 - 14.36 = \qty{89.36}{MPa}$
y $\tau = 28.72\sin\ang{120} = \qty{24.87}{MPa}$.

\textbf{10.} $\sigma_n' = 89.36 - 49.05 = \qty{40.31}{MPa}$.

\textbf{11.} $24.87/40.31 = 0.617$, por debajo de 0.85: la falla aguanta.

\textbf{12.} $0.85\times40.31 = \qty{34.26}{MPa}$; el esfuerzo de cizalla queda
por debajo en \qty{9.39}{MPa}.

\textbf{13.} El círculo tiene su centro en \qty{54.67}{MPa} y radio
\qty{28.72}{MPa}, desde \num{25.95} hasta \qty{83.39}{MPa}; la recta de Byerlee es
$\tau = 0.85\,\sigma_n'$ (todos los \omterm{def:b1:stress-strain-rheology:stress}{esfuerzos normales} aquí están por debajo de
\qty{200}{MPa}); la falla es el punto $(40.31, 24.87)$, a \ang{120} desde el
extremo de $\sigma_1$.

\textbf{14.} $0.85\times54.67/\sqrt{1 + 0.85^2} = 46.47/1.3124 =
\qty{35.41}{MPa}$, más que el radio, \qty{28.72}{MPa}: todo el círculo queda
por debajo de la recta, y ningún plano de la roca puede deslizar.

\textbf{15.} $\phi = \arctan0.85 = \ang{40.36}$: el plano a
$\ang{45} - \phi/2 = \ang{24.8}$ de $\sigma_1$ vertical, con \omterm{def:b1:geological-maps:attitude}{buzamiento}
\ang{65.2}.

\textbf{16.} $S = \sigma_v' = \qty{83.39}{MPa}$ y
$\sigma_1 - \sigma_3 = 2\times0.85\times83.39/(1.3124 + 0.85) =
\qty{65.55}{MPa}$; la pendiente da $13.11\times5 = \qty{65.55}{MPa}$.

\textbf{17.} $\sigma_3 = 132.44 - 65.55 = \qty{66.88}{MPa}$.

\textbf{18.} La falla desliza cuando
$r\sin\ang{120} = 0.85\,(c' + r\cos\ang{120})$ con
$c' = (83.39 + \sigma_3')/2$ y $r = (83.39 - \sigma_3')/2$; resolviendo se obtiene
$\sigma_3' = \qty{17.18}{MPa}$, $\sigma_3 = \qty{66.23}{MPa}$. Es menor que
para el mejor plano: la falla no está en el punto donde el círculo toca
primero la recta, así que el círculo debe crecer más antes de llegar al
punto propio de la falla.

\textbf{19.} Cada \omterm{def:b1:stress-strain-rheology:stress}{esfuerzo normal} efectivo, incluidos $\sigma_1'$ y $\sigma_3'$,
baja en $\Delta P$, y su diferencia no cambia; los esfuerzos de cizalla no
dependen en absoluto de la presión de poros.

\textbf{20.} El mejor plano desliza cuando la distancia iguala al radio:
$0.85\,(54.67 - \Delta P)/1.3124 = 28.72$, así que
$\Delta P = 54.67 - 28.72\times1.3124/0.85 = 54.67 - 44.34 = \qty{10.33}{MPa}$.

\textbf{21.} La columna de agua del pozo pesa lo mismo que antes, así que la
bomba debe añadir los \qty{10.33}{MPa} completos en la boca del pozo, unas 103
veces la presión atmosférica: la altura de una columna de agua
$\num{10.33e6}/(1000\times9.81) = \qty{1053}{m}$.

\textbf{22.} Con $\mu = 0.6$ la falla solo podría soportar
$0.6\times40.31 = \qty{24.19}{MPa}$, menos que los \qty{24.87}{MPa} que
soporta: ya estaría deslizando. Como está quieta, su fricción es al menos
$0.617$, o tiene algo de \omterm{def:b1:stress-strain-rheology:coulomb}{cohesión}.

\textbf{23.} El plano mejor orientado desliza cuando la presión de poros alcanza
$49.05 + 10.33 = \qty{59.38}{MPa}$, es decir $\lambda = 59.38/132.44 = 0.448$,
muy por debajo del valor litostático 1.

\textbf{24.} $\Delta P = \sigma_n' - \tau/0.85 = 40.309 - 24.870/0.85 = 40.309 -
29.259 = \qty{11.049}{MPa}$, más que los \qty{10.33}{MPa} que harían deslizar un
plano de la mejor orientación. Con un decimal, el aumento de presión de poros que
hace deslizar la falla es \textbf{\qty[text-series-to-math=true]{11.0}{MPa}}.
\end{solution}
